\section{Prompt 编写与上下文管理}
%% ============================================================

\subsection{结构化 Prompt 的四要素}

Codex 即使在 prompt 不够完善时也能给出不错的结果，但在大型代码库或高风险任务中，结构化的 prompt 能显著提升可靠性。官方推荐在每个 prompt 中包含以下四个要素：

\begin{importantbox}{Prompt 四要素}
\begin{enumerate}[nosep]
    \item \textbf{Goal}（目标）：你想要修改或构建什么？
    \item \textbf{Context}（上下文）：哪些文件、目录、文档、示例或错误与本任务相关？可用 \texttt{@} 引用文件。
    \item \textbf{Constraints}（约束）：需要遵守哪些标准、架构要求、安全规则或编码规范？
    \item \textbf{Done when}（完成条件）：任务完成时应当满足什么条件？例如测试通过、行为变更或 bug 不再复现。
\end{enumerate}
\end{importantbox}

这四个要素帮助 Codex 限定作用范围、减少假设，并产出更容易 review 的代码。

\subsection{Reasoning Level 选择}

Codex 支持多个推理等级，适配不同任务复杂度：

\begin{table}[H]
\centering
\begin{tabular}{ll}
\toprule
\textbf{推理等级} & \textbf{适用场景} \\
\midrule
Low & 边界清晰的快速任务 \\
Medium / High & 较复杂的代码变更或调试 \\
Extra High & 长链路、高推理需求的 agent 任务 \\
\bottomrule
\end{tabular}
\caption{Reasoning Level 与任务复杂度的对应关系}
\end{table}

\begin{knowledgebox}{语音输入加速上下文提供}
在 Codex App 中，可以使用语音听写（speech dictation）来口述任务需求，而不是手动打字。这对需要提供大量上下文的场景尤其高效。
\end{knowledgebox}

\subsection{复杂任务先规划再执行}

对于复杂、模糊或难以描述的任务，应当让 Codex 先做规划（planning），再进入实现阶段。官方推荐三种方式：

\begin{enumerate}
    \item \textbf{Plan mode}：通过 \texttt{/plan} 或 \texttt{Shift+Tab} 切换。Codex 会先收集上下文、提出澄清问题，再制定实施方案。这是大多数用户最简单有效的选择。
    \item \textbf{让 Codex 采访你}：当你有模糊的想法但不知如何表述时，让 Codex 主动提问、挑战你的假设，把模糊想法转化为具体方案。
    \item \textbf{PLANS.md 模板}：对于更高级的工作流，可配置 Codex 遵循 \texttt{PLANS.md} 或 execution-plan 模板，适用于长期运行或多步骤工作。
\end{enumerate}

\begin{warningbox}{常见错误：跳过规划}
在多步骤和复杂任务中跳过规划（planning）是最常见的错误之一。直接开始编码往往导致方向偏差和大量返工。
\end{warningbox}

\subsection{本章小结}

好的 prompt 包含目标、上下文、约束和完成条件四个要素。根据任务复杂度选择合适的 reasoning level。对于复杂任务，使用 Plan mode 或让 Codex 先采访你，避免直接进入实现。

%% ============================================================
